\section{Directional Shadow and Distance Maps}
\label{sec:directional-shadow-distance-maps}

We precompute directional acceleration maps for maximum intensity projection
(MIP) ray traversal.  The input is a normalized scalar volume
\(V:\mathbb{Z}^{3}\rightarrow[0,1]\), sampled on a regular grid and rendered
with trilinear interpolation.  For MIP, a segment of a ray may be skipped when
all samples inside that segment are guaranteed not to increase the current ray
maximum by more than a prescribed tolerance \(\varepsilon\).  The proposed
maps encode this guarantee as a set of directional binary fields, and then
optionally transform these binary fields into bounded integer skip distances.

\subsection{Ray Classes}

Ray directions are grouped by dominant axis and octant.  A ray class is denoted
by
\[
    g = (a,\sigma),
    \qquad
    a\in\{x,y,z\},
    \qquad
    \sigma\in\{+,-\}^{3},
\]
where \(a\) is the coordinate axis with largest absolute ray component and
\(\sigma\) gives the signs of the ray direction.  The implementation stores
twelve classes: four octant choices for each dominant axis.  Opposite classes
are handled by bidirectional construction, so that the same class represents a
pair of antipodal traversal directions.

For each class \(g\), the volume coordinates are permuted and, if necessary,
reversed so that the computation is performed in a canonical positive
traversal order.  Let \(C_i\) be a cell in this canonical coordinate system,
and let \(F_i^x,F_i^y,F_i^z\) be its three forward faces.  Since the trilinear
interpolant is a convex combination of the eight cell-corner values, extrema
on each face are attained at the four vertices of that face.

\subsection{Extrema-Based Shadow Map}
\label{subsec:extrema-shadow-map}

The first construction explicitly propagates face extrema.  For each cell
\(C_i\) and forward face \(F_i^k\), \(k\in\{x,y,z\}\), we compute the local
face minimum and maximum
\[
    m_i^k = \min_{p\in F_i^k} V(p),
    \qquad
    M_i^k = \max_{p\in F_i^k} V(p).
\]
In practice, the minimization and maximization are over the four vertices of
the face.

The maxima \(M_i^k\) bound the strongest value that the current cell can
contribute through each forward face.  The minima are propagated through the
volume to form a conservative lower bound on the maximum value that any
monotone ray prefix must have already encountered before entering the cell.
Let \(L_i^k\) denote this propagated lower bound for face \(F_i^k\).  The
propagation has the form
\[
    L_i^k =
    \max\left(
        m_i^k,\,
        \min_{(j,\ell)\in\mathcal{P}(i,k)} L_j^\ell
    \right),
    \qquad k\in\{x,y,z\},
\]
where \(\mathcal{P}(i,k)\) is the set of predecessor faces that can reach
\(F_i^k\) under the monotone traversal associated with \(g\).  The minimum over
predecessors makes the test conservative across all admissible monotone paths,
while the outer maximum incorporates the current face values into the running
MIP lower bound.

A cell is marked as shadowed for class \(g\) when its possible local maximum is
already bounded by the propagated prefix value:
\[
    S_g(i)=1
    \quad\Longleftrightarrow\quad
    M_i^k - L_i^k < \varepsilon
    \quad\forall k\in\{x,y,z\}.
\]
Thus, if \(S_g(i)=1\), no ray of class \(g\) can obtain a MIP update larger
than \(\varepsilon\) by evaluating \(C_i\).  The cell may therefore be skipped
without changing the image beyond the selected tolerance.

\paragraph{Bidirectional construction.}
For each stored ray class, a forward shadow map is computed first.  A reverse
pass is then evaluated in the complementary direction.  During the reverse
pass, cells that were already classified as shadowed by the forward pass are
treated as hollow, so that previously skipped regions do not become artificial
occluders for the opposite direction.  The final bidirectional shadow field is
the logical union
\[
    S_g^{\mathrm{bi}}(i)
    =
    S_g^{\mathrm{fwd}}(i)
    \lor
    S_g^{\mathrm{bwd}}(i).
\]

\subsection{Difference-Based Shadow Map}
\label{subsec:difference-shadow-map}

The second construction computes the same kind of directional shadow
classification in a difference domain, avoiding explicit storage of separate
minimum and maximum maps.  For a canonical traversal direction, each grid
position stores four differences between predecessor samples on the incoming
face and the current sample:
\[
    \Delta_r(i) = V(i-r)-V(i),
    \qquad r\in\mathcal{R},
\]
where \(\mathcal{R}\) indexes the four incoming-face offsets.  A positive
difference means that the predecessor already contains an intensity at least
as large as the current sample.

The propagation accumulates the guaranteed positive margin available from
predecessor states.  For a predecessor \(i-r\), let
\[
    q(i-r)=\min_{s\in\mathcal{R}}\Delta_s(i-r)
\]
be the worst guaranteed margin at that predecessor.  The propagated update is
\[
    \Delta_r(i)
    \leftarrow
    \Delta_r(i) + \max(q(i-r),0).
\]
This recurrence is the difference-domain analogue of propagating a lower
bound on the already-seen ray maximum.  It accumulates only the part of the
predecessor guarantee that is safely non-negative.

After propagation, the current cell is shadowed when all relevant outgoing
vertex differences remain non-negative within tolerance:
\[
    S_g(i)=1
    \quad\Longleftrightarrow\quad
    \min_{v\in F_i^+}\min_{r\in\mathcal{R}}\Delta_r(v)
    >
    -\varepsilon.
\]
Equivalently, every sample that the current cell could contribute is already
dominated by an upstream sample, up to the tolerance \(\varepsilon\).

The bidirectional difference-based method follows the same structure as the
extrema-based method.  First, a forward shadow map is computed.  Then a set of
gates is derived from this forward map and used in the reverse propagation.  A
closed gate clamps the recurrence to positive margins, while an open gate lets
the signed margin pass through.  This prevents the reverse pass from using
regions already eliminated by the forward pass as invalid occluders.  The
forward and reverse binary maps are finally merged with a logical union.

\subsection{Block Shadow Maps}

The directional maps may be reduced to block resolution.  Let \(B_b\) be the
set of cells contained in block \(b\).  The block-level shadow value is
computed by min-pooling the binary cell map:
\[
    \bar{S}_g(b)
    =
    \min_{i\in B_b} S_g(i).
\]
This operation is conservative: a block is marked shadowed only if every cell
inside it is shadowed for the corresponding ray class.  The resulting block
map has dimensions
\[
    \left\lceil \frac{N_x+1}{h} \right\rceil
    \times
    \left\lceil \frac{N_y+1}{h} \right\rceil
    \times
    \left\lceil \frac{N_z+1}{h} \right\rceil,
\]
where \(h\) is the block size and \(N_x,N_y,N_z\) are the volume dimensions.

\subsection{Extended Anisotropic Directional Map}

The complete acceleration representation consists of twelve directional
fields.  We store four classes for dominant \(x\), four for dominant \(y\), and
four for dominant \(z\):
\[
    \mathcal{G}
    =
    \mathcal{G}_x
    \cup
    \mathcal{G}_y
    \cup
    \mathcal{G}_z,
    \qquad
    |\mathcal{G}|=12.
\]
At render time, the ray dominant axis and octant select the corresponding
field.  The maps can be bit-packed into compact integer textures, so one
texture lookup returns the shadow information for all relevant directional
groups at a block.

\subsection{Shadow Distance Maps}
\label{subsec:shadow-distance-maps}

The binary shadow maps are further transformed into bounded distance maps.  A
binary field only answers whether the current block can be skipped; a distance
field estimates how many consecutive blocks can be skipped safely.

For a directional block shadow map \(\bar{S}_g\), the initial distance field is
\[
    D_g^{0}(b)
    =
    \begin{cases}
        R, & \bar{S}_g(b)=1,\\
        0, & \bar{S}_g(b)=0,
    \end{cases}
\]
where \(R\) is the maximum representable distance for the chosen texture
format.  Non-shadowed blocks are obstacles with distance zero, and shadowed
blocks are initialized to the largest possible skip length.

\paragraph{Isotropic Chebyshev distance.}
The isotropic variant computes a capped Chebyshev distance to the nearest
non-shadowed block:
\[
    D_g(b)
    =
    \min\left(
        R,\,
        \min_{c:\bar{S}_g(c)=0}
        \|b-c\|_{\infty}
    \right).
\]
This is implemented by separable one-dimensional dynamic passes along
\(x\), \(y\), and \(z\).  Along an axis \(a\), the update is
\[
    D'(b)
    =
    \min_{n}
    \max\left(D(b+n e_a),\, |n|\right),
\]
with the search truncated by \(R\) and by the volume boundary.

\paragraph{Extended anisotropic distance.}
The anisotropic variant restricts the distance computation to the oriented
cone associated with the current ray class.  Let the two non-dominant axes be
\(u\) and \(v\), and let the dominant axis be \(a\).  The algorithm performs
one-sided Chebyshev passes along \(u\) and \(v\) in the octant signs of
\(\sigma\), followed by an extended one-sided pass along \(a\):
\[
    D^{u}(b)
    =
    \min_{n\geq 0}
    \max\left(D^{0}(b+n\sigma_u e_u),\, n\right),
\]
\[
    D^{v}(b)
    =
    \min_{n\geq 0}
    \max\left(D^{u}(b+n\sigma_v e_v),\, n\right),
\]
and
\[
    D_g(b)
    =
    \min\left\{
        n\geq 1:
        D^{v}(b+n\sigma_a e_a) \leq n
    \right\},
\]
clamped to \(R\) if no such \(n\) is found.  The first two passes enlarge the
valid skip region transversely, while the final dominant-axis pass determines
how far the ray may advance before it can encounter a non-shadowed block.

A bidirectional anisotropic distance can also be formed by computing the same
oriented distance in the opposite sign direction and taking the minimum:
\[
    D_g^{\mathrm{bi}}(b)
    =
    \min\left(D_g^{+}(b),D_g^{-}(b)\right).
\]

\subsection{Texture Packing And Runtime Use}

The distance fields are uploaded as integer 3D textures.  Four storage variants
are used:
\[
\begin{array}{c|c|c}
\text{variant} & \text{texture format} & \text{stored information} \\
\hline
\text{1-bit}  & \mathrm{R16UI}    & 12 binary shadow flags \\
\text{5-bit}  & \mathrm{RGBA16UI} & x,y,z distances packed per quadrant \\
\text{8-bit}  & \mathrm{RGB32UI}  & four 8-bit distances per axis \\
\text{10-bit} & \mathrm{RGBA32UI} & larger packed distances per quadrant
\end{array}
\]
At runtime, the fragment shader determines the ray group \(g\), fetches the
corresponding packed value, unpacks the skip distance \(D_g(b)\), and sets
\[
    \mathrm{empty}(b) = [D_g(b)>0].
\]
If \(D_g(b)=0\), the block is traversed normally.  If \(D_g(b)>0\), the ray is
advanced across a block region of radius \(D_g(b)\).  Therefore, the distance
map acts as a texture-resident wrapper around the directional shadow maps:
the original binary visibility guarantee is preserved, while adjacent shadowed
regions are compressed into larger skip operations.

